\documentclass[10pt,a4paper]{amsart}
\usepackage[margin=19mm]{geometry}
\usepackage{amsmath,amssymb,mathtools}
\usepackage[scr=rsfs,cal=euler]{mathalfa}
\usepackage{xfrac}
\usepackage[unicode,hidelinks,bookmarks=false]{hyperref}
\newcommand{\ie}{i.e.\ }
\newcommand{\cf}{cf.\ }
\newcommand{\resp}{respectively\ }
\newcommand{\Subsection}{\subsection}
\providecommand{\nobreakdash}{\nobreak\hbox{-}}
\setlength{\emergencystretch}{2em}
\multlinegap0pt
\usepackage{amssymb}
\usepackage{xcolor}
\usepackage{csquotes}
\usepackage{enumerate}
\usepackage[normalem]{ulem}
\usepackage{cancel}
\newtheorem{lemma}{Lemma}
\newcommand{\R}{\mathbb{R}}
\newcommand{\T} {\mathbb T}
\newcommand{\bq}{\begin{equation}}
\newcommand{\calJ}{\mathcal J}
\newcommand{\dx}{\textnormal{d}x}
\newcommand{\eq}{\end{equation}}
\newcommand{\intt}{\int_{\T^d}}
\newcommand{\lt}{\left}
\newcommand{\rt}{\right}
\title{Global existence of Lagrangian solutions to the ionic Vlasov--Poisson system}
\author{Young-Pil Choi, Dowan Koo, Sihyun Song}
\date{Source: arXiv:2501.13872v2 (CC BY 4.0)}
\begin{document}
\maketitle

\noindent\emph{Source and licence.} Global existence of Lagrangian solutions to the ionic Vlasov--Poisson system. Version arXiv:2501.13872v2 (CC BY 4.0), distributed by arXiv under the Creative Commons Attribution 4.0 International licence, http://creativecommons.org/licenses/by/4.0/, as recorded in the arXiv licence metadata for this exact version and shown on the abstract page https://arxiv.org/abs/2501.13872v2. Full licence notice: \texttt{licenses/arxiv-2501.13872-CC-BY-4.0.txt}.\par\smallskip

\noindent\textit{Editorial scope.} Counted unit: the article's lemma lem:C1 (source0.tex lines 920--925). The article's complete original proof of it is delivered with it (source0.tex lines 926--970, one \texttt{proof} environment of 45 lines). No mathematics is written, repaired or re-derived: the statement and the proof are the article's own lines, in the article's own notation, and no retained line is substituted or re-flowed.  Retained uncounted as the source-local context the proof needs: lines 781--786; lines 789--814; lines 797--801.  Nothing was omitted from any retained span, so the package carries no omission record.  No retained line contains a citation, so the wrapper carries no bibliography and no bibliography entry.  No retained line contains a figure environment, an image-inclusion command or drawing code, so no image is delivered and the package declares no asset dependency.  Editorial formatting only, in addition: an AMS \texttt{amsart} preamble that restates the article's own declarations and macros used by the retained lines, namely the environments \texttt{lemma} on separate counters, together with the standard packages those lines need.  The retained environments print in this document as Lemma 1 in order of appearance, whereas the article prints the same items with the numbering of its own full document.  The complete native source of the article as distributed in the arXiv e-print for this exact version is bundled unchanged as \texttt{sources/172058/source0.tex}.
\begin{equation}\label{eq:jfm}
\begin{split}
e^{\Phi_\flat + h} - \lt(m_\flat + \rho^\sharp \rt)h &=  e^{\Phi_\flat}e^{h}-\lt(m_\flat + \rho^\sharp \rt) h_+ + \lt(m_\flat + \rho^\sharp \rt) h_- \\
&\ge e^{-\|\Phi_\flat\|_{L^\infty}}e^{h}-m_\flat h_+ - \frac{(h_+)^{p'}}{p'}- \frac{\rho^p}{p} + \lt(m_\flat + \rho^\sharp \rt) h_- ,
\end{split}
\end{equation}

When $h<0$, it is clear that \eqref{eq:jfm} can be estimated further as
\begin{equation}
\label{eq: h<0}
\begin{split}
    e^{\Phi_\flat + h} - \lt(m_\flat + \rho^\sharp \rt)h & \ge e^{-\|\Phi_\flat\|_{L^\infty}}e^h + \lt(m_\flat + \rho^\sharp\rt) h_- - \frac{\rho^p}{p}  \ge \frac{1}{2}e^{-\|\Phi_\flat\|_{L^\infty}}e^h + m_\flat |h| + \rho^\sharp h_- - \frac{\rho^p}{p}.
    \end{split}
\end{equation}
On the other hand, when $h\ge 0$ we use the convexity of $x \mapsto e^x$ and its Taylor expansion to yield
\begin{equation}\label{eq:c1}
\begin{split}
e^{-\|\Phi_\flat\|_{L^\infty}}e^{h}-m_\flat h_+ - \frac{(h_+)^{p'}}{p'} &\ge  \frac{1}{2}e^{-\|\Phi_\flat\|_{L^\infty}}e^{h} + m_\flat h_+ -C_1 = \frac{1}{2}e^{-\|\Phi_\flat\|_{L^\infty}}e^h + m_\flat |h| - C_1,
\end{split}
\end{equation}
where 
\bq\label{eq:c1f}
C_1:= N_d(\|\rho\|_{L^1(\T^d)} + \|\rho\|_{L^1(\T^d)}^2 ) + (4e^{N_d\|\rho\|_{L^1(\T^d)}} \Gamma(p'+3))^{p'} >0
\eq
(see Lemma \ref{lem:C1} below for details). Consequently, for $h\ge 0$, \eqref{eq:jfm} can be estimated further as
\begin{align*}
    e^{\Phi_\flat+h} - \lt(m_\flat + \rho^\sharp \rt)h \ge \frac{1}{2}e^{-\|\Phi_\flat\|_{L^\infty}} e^h + m_\flat |h| - C_1 - \frac{\rho^p}{p}.
\end{align*}
Together with \eqref{eq: h<0}, we deduce that
\bq\label{eq:jlb}
\calJ[h] \ge \intt \frac{1}{2}|\nabla h|^2 + \frac{1}{2}e^{-\|\Phi_\flat\|_{L^\infty}}e^{h} +m_\flat |h| + \rho^\sharp h_-\,\dx -C_1 - \frac{1}{p}\intt \rho^p \,\dx
\eq
which confirms that $\calJ[\cdot]$ is bounded below.

\begin{equation}
\begin{split}
e^{-\|\Phi_\flat\|_{L^\infty}}e^{h}-m_\flat h_+ - \frac{(h_+)^{p'}}{p'} &\ge  \frac{1}{2}e^{-\|\Phi_\flat\|_{L^\infty}}e^{h} + m_\flat h_+ -C_1 = \frac{1}{2}e^{-\|\Phi_\flat\|_{L^\infty}}e^h + m_\flat |h| - C_1,
\end{split}
\end{equation}

\begin{lemma}\label{lem:C1}
The constant $C_1$ in \eqref{eq:c1f} can be estimated as
\[
C_1 \le N_d(\|\rho\|_{L^1(\T^d)} + \|\rho\|_{L^1(\T^d)}^2 ) + (4e^{N_d\|\rho\|_{L^1(\T^d)}} \Gamma(p'+3))^{p'}.
\]
\end{lemma}

\begin{proof}
We aim to obtain the estimate for the constant $C_1>0$ in \eqref{eq:c1}, which satisfies
\begin{align*}
e^{-\|\Phi_\flat\|_{L^\infty}}e^{h}-m_\flat h - \frac{h^{p'}}{p'} &\ge  \frac{1}{2}e^{-\|\Phi_\flat\|_{L^\infty}}e^{h} + m_\flat h -C_1, \quad \forall \, h \geq 0.
\end{align*}


From the  convexity and positivity of $x \mapsto e^x$, we get
\[
e^h \ge e^s + e^s(h-s) \ge e^s(h-s) \quad \forall \, h,s\in \R.
\]
By choosing $s:=a +\log4 + \log(2b)$, where $a,b>0$, we find
\[
\frac{1}{4}e^{-a}e^h - bh \ge bh -(2b)(a+\log4 +\log(2b)).
\]
Now, choose $a:=\|\Phi_\flat\|_{L^\infty(\T^d)}\le N_d\|\rho\|_{L^1(\T^d)}$ and $b :=m_\flat \le \|\rho\|_{L^1(\T^d)}$. Then we obtain
\begin{equation}
\label{eq: poly}
\frac{1}{4}e^{-\|\Phi_\flat\|_{L^\infty}}e^h - m_\flat h \ge m_\flat h - N_d(\|\rho\|_{L^1(\T^d)} +\|\rho\|_{L^1(\T^d)}^2)
\end{equation}
for some constant $N_d>0$. On the other hand, we also note that
\begin{align}
\label{eq : h >>}
\frac{1}{4}e^{-a+h} - \frac{h^{p'}}{p'} \ge 0,\quad \forall \, h \ge 4e^a m!>1,\quad m:= \left\lceil p' \right\rceil +1.
\end{align}
Indeed, for such $h$, we deduce
\[
\frac{1}{4}e^{-a+h} = \frac{e^h}{4e^a} \ge \frac{h}{4e^am!}h^{m-1} \ge h^{p'} \ge \frac{h^{p'}}{p'}.
\]
In the opposite case where $0\le h< 4e^a m!$, we can simply estimate
\begin{equation*}
\begin{split}
\frac{1}{4}e^{-a+h} - \frac{h^{p'}}{p'}  \ge -\frac{h^{p'}}{p'}  \ge -(4e^a m!)^{p'} \ge - (4e^a\Gamma(p'+3))^{p'}.
\end{split}
\end{equation*}
This and \eqref{eq : h >>} imply
\begin{align*}
    \frac{1}{4}e^{-a+h} - \frac{h^{p'}}{p'} \ge -(4e^a \Gamma(p' + 3))^{p'} \quad \forall \, h\ge 0.
\end{align*}
Adding this with \eqref{eq: poly}, we obtain
\begin{align*}
    \frac{1}{2}e^{-a+h} - m_\flat h - \frac{h^{p'}}{p'} \ge m_\flat h - N_d(\|\rho\|_{L^1(\T^d)} + \|\rho\|_{L^1(\T^d)}^2 ) - (4e^{N_d\|\rho\|_{L^1(\T^d)}} \Gamma(p'+3))^{p'}.
\end{align*}
This completes the proof.
\end{proof}

\end{document}
